\scanpage{139}
\mspage{55}

\[
M\boxtimes N
\tag*{(4.1.9)}
\]

par la formule

\deleted{\[
\begin{aligned}
M\boxtimes N &\unsure{\simeq|=} \text{sous-catégorie pleine de}\\
&\underline{\operatorname{Hom}}^{!}(M^{\mathrm{o}},N)
\end{aligned}
\tag*{(\unsure{4.1.10})}
\]}

\[
M\boxtimes N\simeq\underline{\operatorname{Hom}}^{!!}(M^{\mathrm{o}}\times N^{\mathrm{o}},\mathrm{Ens})
\tag*{(4.1.10)}
\]

où le signe $\underline{\operatorname{Hom}}^{!!}$ désigne la sous-catégorie pleine \unsure{dans|de} la catégorie $\underline{\operatorname{Hom}}$, formée des bi-contrafoncteurs \deleted{\unsure{$F$|$B$}}

\[
F:M^{\mathrm{o}}\times N^{\mathrm{o}}\unsure{\longrightarrow|\text{---}}\mathrm{Ens}
\tag*{(4.1.10.1)}
\]

tels que pour tout $x$ dans $M$, le foncteur

\[
y\longmapsto F(\unsure{x|a},y):N^{\mathrm{o}}\unsure{\text{---}|\longrightarrow}\mathrm{Ens}
\tag*{(4.1.10.2)}
\]

\deleted{\unsure{transforme}} soit dans $\underline{\operatorname{Hom}}^{!}(N^{\mathrm{o}},\mathrm{Ens})$. i.e transforme \inserted{petites} $\underset{\longrightarrow}{\lim}$ en $\underset{\longleftarrow}{\lim}$, et que, symétriquement, pour \unsure{tt|tout} $y$ dans $N$, le foncteur

\[
x\longmapsto F(x,y):M^{\mathrm{o}}\text{---}\mathrm{Ens}
\tag*{(4.1.10.3)}
\]

soit dans $\underline{\operatorname{Hom}}^{!}(M^{\unsure{\mathrm{o}|b}},\mathrm{Ens})$. Mais cela signifie aussi que ces foncteurs sont \uline{représentables} ; donc on trouve aussi

\[
M\boxtimes N\xleftarrow{\approx}\underline{\operatorname{Hom}}^{!}(N^{\mathrm{o}},M)
\tag*{(4.1.11)}
\]

et

\[
M\boxtimes N\xleftarrow{\approx}\underline{\operatorname{Hom}}^{!}(M^{\mathrm{o}},N)
\tag*{(4.1.12)}
\]
% archive sheet number: 138
